\chapter{VAMP --- Formal Definition and Properties}\label{chap:vamp}

[todo: probably wanna redo this para to follow nicely from prev chapter on how we need a game theoretic formulation of decentralized ITAP. and then whhy markets are needed idk]


\renewcommand{\proves}{\vdash}
We formalize a \textbf{Verifiable Allocation Market POSG (VAMP)} as a market-mediated, partially observable stochastic game that arises when HTAP is decentralized: agents allocate compute across proof/conjecture/publication actions and may additionally trade \textbf{verifiable, state-contingent contracts} whose settlement is triggered by a public \textbf{Verify} oracle. This section gives a LaTeX-ready VAMP tuple aligned with our existing library/market notation and two-stage transition construction. We then define a contract space (bounties, securities/claims, bundles), admissibility/collateral rules, and a modular market mechanism interface \(\mathcal{M}\). We instantiate canonical mechanisms---posted bounties, bilateral OTC contracts, convex-cost/LMSR and CFMM-style AMMs, and combinatorial auctions---and state the appropriate equilibrium concepts (Bayesian Nash, subgame-perfect refinements, Markov perfect equilibria) with existence remarks grounded in fixed-point theorems. Finally, we analyze incentive alignment (bounties internalize publication externalities via a short proof), manipulation modes (front-running, sybil, withholding), core complexity barriers (PPAD/NEXP), and baseline algorithmic benchmarks that connect directly to the implementation and MARL chapters.


\subsection{Partially Observable Stochastic Games}


A \emph{partially observable stochastic game} (POSG) is a mathematical model for sequential decision-making with multiple interacting agents under partial information. In this work we focus on the standard \emph{infinite-horizon discounted} formulation.
\subsubsection{Definition}

An infinite-horizon discounted POSG is specified by the tuple
\[
\Big(\mathcal{N}, \mathcal{W},\{\mathcal{O}^\alpha\}_{\alpha\in \mathcal{N}},\{\mathcal{Z}^\alpha\}_{\alpha\in \mathcal{N}},
\{\mathcal{A}^\alpha\}_{\alpha\in \mathcal{N}},\]\[\{\mathrm{Adm}^\alpha\}_{\alpha\in\mathcal{N}},\mathcal{T},
\{r^\alpha\}_{\alpha\in \mathcal{N}},\gamma,b_0\Big)    
\]

where:
\begin{itemize}
    \item $\mathcal{N}$ is the set of agents, indexed by $\alpha\in \mathcal{N}$.
    \item $\mathcal{W}$ is the world (state) space.
    \item $\mathcal{A}^\alpha$ is the (ambient) action space of agent $\alpha$, and the joint action space is
    $\mathcal{A}:=\prod_{\alpha\in \mathcal{N}}\mathcal{A}^\alpha$.
    \item $\mathcal{O}^\alpha$ is the observation space of agent $\alpha$.
    \item $\mathrm{Adm}^\alpha:\mathcal{W}\times \mathcal{A}^\alpha\to\{0,1\}$ is the \emph{admissibility predicate} for agent $\alpha$.
    It induces the state-dependent legal action set
    \[
    \mathcal{A}^\alpha(s) := \{a^\alpha\in\mathcal{A}^\alpha : \mathrm{Adm}^\alpha(s,a^\alpha)=1\}.
    \]
    We assume $\mathcal{A}^\alpha(s)\neq\varnothing$ for all $s$ (e.g.\ containing a no-op action).
    \item $\mathcal{T}$ is the transition kernel: for $s,s'\in \mathcal{W}$ and joint action $a\in\mathcal{A}$,
    \[
    \mathcal{T}(s' \mid s,a)\in[0,1]
    \]
    is the probability of transitioning from $s$ to $s'$ under $a$.
    \item $\mathcal{Z}^\alpha$ is the observation kernel for agent $\alpha$: for $o^\alpha\in\mathcal{O}^\alpha$,
    \[
    \mathcal{Z}^\alpha(o^\alpha \mid s,a,s')\in[0,1]
    \]
    is the probability that agent $\alpha$ observes $o^\alpha$ given transition $s\xrightarrow{a}s'$.
    \item $r^\alpha:\mathcal{W}\times \mathcal{A}\to \mathbb{R}$ is the per-step reward function for agent $\alpha$.
    \item $\gamma\in[0,1)$ is the discount factor.
    \item $b_0$ is the initial distribution over states $s_0\sim b_0$.
\end{itemize}

At each timestep $t$, each agent selects an action $a_t^\alpha\in \mathcal{A}^\alpha(s_t)$, yielding a joint action
$a_t=(a_t^\alpha)_{\alpha\in\mathcal{N}}\in\mathcal{A}$.
(Equivalently, one may allow agents to propose any $a^\alpha\in\mathcal{A}^\alpha$ and define a convention for inadmissible actions,
such as rejection/no-op or penalties; we use explicit admissibility predicates for clarity.)

\subsubsection{Policies, Utilities, and Equilibria}

A (behavioral) policy for agent $\alpha$ is a mapping from the player's action--observation history to a distribution over actions.
Let
\[
h^\alpha_t := (o^\alpha_0,a^\alpha_0,o^\alpha_1,a^\alpha_1,\dots,o^\alpha_t)\in(\mathcal{O}^\alpha\times \mathcal{A}^\alpha)^t\times \mathcal{O}^\alpha
\]
denote agent $\alpha$'s history up to time $t$. A policy $\pi^\alpha(\cdot \mid h^\alpha_t)$ is a distribution over $\mathcal{A}^\alpha$ conditioned on the past history $h^\alpha_t$. We write $\Pi^\alpha$ for the space of policies available to agent $\alpha$.

A \emph{pure strategy} is a map $\sigma^\alpha$ that selects a single action deterministically as a function of $h^\alpha_t$. A joint policy is $\pi := (\pi^\alpha)_{\alpha\in \mathcal{N}}$, and we write $\pi^{-\alpha}$ for the policies of all players except $\alpha$.

Given a joint policy $\pi$, player $\alpha$'s expected discounted return is
\[
R^\alpha(\pi) := \mathbb{E}_{\mathcal{T},\mathcal{Z},\pi}\!\left[\sum_{t=0}^{\infty}\gamma^t\, r^\alpha(s_t,a_t)\;\middle|\; s_0\sim b_0\right]
\]
where the expectation is taken over the the transition kernels $T$, observation kernels $Z$, and the (possibly stochastic) action selection induced by $\pi$. Under bounded rewards (e.g.\ $|r^\alpha(s,a)|\le R_{\max}$) and $\gamma\in[0,1)$, this sum converges almost surely
and $|R^\alpha(\pi)|\le \frac{R_{\max}}{1-\gamma}$.

Because $R^\alpha(\pi)$ depends on the joint policy and reward functions need not align across players, a POSG is not generally a single-agent optimization problem. Rather, one studies various behavioral policies and game-theoretic equilibria. One such equilibrium is the \emph{Nash equilibrium}: a joint policy $\pi$ is a Nash equilibrium if every agent $\alpha$ plays a best response to $\pi^{-\alpha}$, i.e., for all policies $\pi^{\alpha'}\in \Pi^\alpha$,
\[
R^\alpha(\pi^\alpha,\pi^{-\alpha}) \;\ge\; R^{\alpha}(\pi^{\alpha'},\pi^{-\alpha})
\]

% This is a rather strong condition, and one may also want to study weaker notions of equilibrium behaviors, such as correlated equilibria or coarse-correlated equilibria. \textbf{TODO DEFINE!}

A particularly important taxonomic feature of POSGs is the relationship between agents' reward functions. In \emph{cooperative} games, all agents share the same reward function: $r^\alpha(s,a)=r^{\alpha'}(s,a)$ for all $i,j$. In a two-player \emph{zero-sum} game,
$r^0(s,a)=-r^1(s,a)$ for all $(s,a)$. When no such restriction is imposed, the POSG is termed \emph{general-sum}.


\subsection{VAMP Preliminaries}


\subsubsection{Syntax}

We consider the \textit{formula} $\phi$ as the core syntactic primitive of (formal) mathematics. Namely, we represent mathematical theorem statements syntactically as a formula (i.e. object of type \texttt{Prop}) within some formal language $\Sigma$, such as Lean, Rocq, etc. We define $\mathcal{F}$ as the collection of all such formulas, and equip it with an involution $\lnot : \mathcal{F}\rightarrow \mathcal{F}$; namely, $\lnot(\lnot\varphi)=\varphi$ for all $\varphi\in\mathcal{F}$. Note that we will refer to formulas as ``nodes'' interchangably due to a natural graphical structure that will be discussed later.

\subsubsection{Semantics}

We now aim to represent the semantic meaning of a theorem statement within a particular context concisely. Towards this, we consider the \textit{sequent}, given as an element of $\mathcal{S} := \mathcal{P}_{fin}(\mathcal{F}) \times \mathcal{F} $ and stylized as $[\Gamma \proves \phi]$ (read: ``Gamma \textit{proves} phi''). Namely, such a sequent represents a \textit{proof obligation}, i.e. the notion of proving a theorem statement $\phi$ (the \textit{consequent}) within the \textit{antecedent} context $\Gamma$. Such a proof obligation may be \textit{resolved}, meaning proven.

\subsubsection{Libraries}

We assume the existence of a universal clock $t \in \mathbb N$ and set of agents $\mathcal N$, and use it to define a library $\mathcal{L}_\bullet = \{\mathcal{L}_t\}_{t \in \mathbb{N}}$ as a collection of discrete library states that evolves over time. At every discrete timestep $t \in \mathbb{N}$, each library state is given as a tuple 
\[\mathcal{L}_t := (\mathcal{C}_t,\mathcal{RS}_t, \delta_t, \Theta_t) \in \mathscr{L}\]
where we partition $\mathcal{F} = \mathcal{C}_t \cup \mathcal{G}_t, \mathcal{C}_t \cap \mathcal{G}_t = \varnothing$ into \textit{concrete} and \textit{ghost} nodes. Intuitively, concrete nodes represent theorem statements that have been stated or conjectured, and ghost nodes represent all hypothetical, possible theorem statements that are yet to be conjectured or stated. $\mathcal{C}_t$ naturally induces a semantic analogue: let $\mathcal{CS}_t := \mathcal{P}_{fin}(\mathcal{C}_t) \times \mathcal{C}_t$ be the set of all \textit{concrete sequents}, which contains a subset $\mathcal{RS}_t \subseteq \mathcal{CS}_t$ of \textit{resolved sequents}, which represent the concrete nodes that have been proven (within a particular antecedent context, of course). 


Mathematically, a proof of a particular theorem in some context often explicitly depends on other lemmas within this context. In an effort to model this phenomenon, define $\delta_t : \mathcal{RS}_t \to \mathcal{P}_{fin}(\mathcal{F})$ as the dependency mapping.
Note that in mathematics, there may exist many proofs of the same theorem with different dependencies for each, and one may find new proofs of the same theorem that are also valid. Thus, in the most general formulation, we only require that resolved sequents must remain resolved ($\mathcal{RS}_{t} \subseteq \mathcal{RS}_{t+1}$), though the underlying ``proof'' (moreso, the set of dependencies of this ``proof'') can change.

And lastly, we define $\Theta_t : \mathcal{RS}_t\to\mathbb N^+ \times \mathcal N$ to record the (discrete) solve-time associated to each resolved sequent, and the agent responsible for the resolution. 
% as in our formulation we do not allow the modification or deletion of a proof after it is initially proven (as $\mathcal{RS}_{t_0} \subseteq \mathcal{RS}_{t}$ for $t_0 < t$).


So, in other words, the library state $\mathcal{L}_t$ represents the collection of all stated theorem statements, the collection of proof witnesses within particular contexts for a subset of these statements, and the dependencies and ancillary metadata of these witnesses.

\paragraph{Library Invariants and Properties}

Naturally, we impose a monotonicity invariant for libraries: one may not ``unconjecture'' or ``unprove'' a theorem. Formally, given a library $\mathcal{L}_\bullet = \{\mathcal{L}_t\}_{t \in \mathbb{N}}=\{(\mathcal{C}_t,\mathcal{RS}_t,\delta_t,\Theta_t)\}_{t \in \mathbb{N}}$, we require that for all $t \in \mathbb{N}$:
\[\mathcal{C}_{t} \subseteq \mathcal{C}_{t+1} \quad \text{and} \quad \mathcal{RS}_{t} \subseteq \mathcal{RS}_{t+1}\]

Additionally, we require that if one conjectures a theorem, they must also have automatically conjectured its converse; in other words:
\[\phi \in \mathcal{C}_t \iff \neg\phi \in \mathcal{C}_t\]

Similarly, one may not prove both a theorem and its converse true: 
\[[\Gamma \proves \phi] \in \mathcal{RS}_t \implies [\Gamma \proves \neg \phi] \notin \mathcal{RS}_t\]

Semantically, we note reflexivity holds: $\phi \in \Gamma \implies [\Gamma \proves \phi] \in \mathcal{RS}_t$. And for the sake of consistency, we require that for any $[\Gamma \proves \phi] \in \mathcal{CS}_t$, $\psi \in \Gamma \implies \neg\psi \notin \Gamma$ (although we do not impose any full propositional semantics or richer logical theory behind these abstract formulas, preferring to leave them fully symbolic objects with various resolution and concreteness states, allowing inconsistent antecedents in our witnesses limits the scalability of our approach to deployment as a strategy module for multiagent theorem proving in a fully-defined formal language and type theory). 

Note that this last condition implicitly redefines
\[\mathcal{CS}_t := \mathcal{P}_{fin}(\mathcal{C}_t)/\sim \times \mathcal{C}_t \] where $\phi\sim\neg\phi$ is the equivalence relation generated by the negation operation. So, due to this consistency condition, we note that $\mathcal{RS}_t$ is upwards-closed, in that if $[\Gamma \proves \phi] \in \mathcal{RS}_t$, then for all $\Gamma' \in \mathcal{P}_{fin}(\mathcal{C}_t)$, $\Gamma \subseteq \Gamma' \implies [\Gamma'\proves\phi]\in \mathcal{RS}_t$.

Additionally, considering the dependency mapping, we note that intuitively $\delta_t$ must satisfy the property that $\delta_t([\Gamma \proves \phi]) \subseteq \Gamma$. Moreover, one may use $\delta_t$ to induce a directed graph $D(\delta_t) := (\mathcal{CS}_t,\mathcal{E}_t)$, where 
\[\mathcal{E}_t = \Big\{([\Gamma\proves \phi], [\Gamma' \proves \psi]) : \phi \in \delta_t([\Gamma' \proves \psi]), \Gamma\subseteq \Gamma'\Big\}\]
We require that this induced directed graph $D(\delta_t)$ of $\delta_t$ must be acyclic.

\paragraph{Axioms and Fully Resolved Sequents}

As an aside, we additionally note that although $[\Gamma \proves \phi] \in \mathcal{RS}_t$ may be marked resolved, that does not mean that $\phi$ has a ``error''-free proof within the current library state as its proof may depend on some still-unresolved lemmas. More concretely, we may define $\mathcal{FRS}_t \subseteq \mathcal{RS}_t$ as the collection of \textit{fully resolved sequents}, which is defined inductively such that \textit{axioms} $[\proves \phi] \in \mathcal{FRS}_t$ and $[\Gamma \proves \phi] \in \mathcal{FRS}_t$ if and only if for all $\phi' \in \delta_t([\Gamma \proves \phi])$, there exists $\Gamma'\subseteq \Gamma$ such that $[\Gamma' \proves \phi'] \in \mathcal{FRS}_t$.


\paragraph{Private and Shared Libraries}
In the multiagent setting, each agent $\alpha\in\mathcal{N}$ maintains a \textit{private library} $\mathcal{L}_\bullet^\alpha = \{\mathcal{L}^\alpha_t\}_{t\in\mathbb{N}}$ with
\[
\mathcal{L}^\alpha_t := (\mathcal{C}^\alpha_t,\mathcal{RS}^\alpha_t,\delta^\alpha_t,\Theta_t^\alpha),
\]
and the environment maintains a \textit{shared (public) library} $\mathcal{L}_\bullet^0 = \{\mathcal{L}^0_t\}_{t\in\mathbb{N}}$ with
\[
\mathcal{L}^0_t := (\mathcal{C}^0_t,\mathcal{RS}^0_t,\delta^0_t, \Theta_t^0).
\]

We additionally enforce the following inclusion invariant for all $t \in \mathbb{N}, \alpha \in \mathcal{N}$:
\[
\mathcal{L}^0_t \subseteq \mathcal{L}^\alpha_t
\]
meaning $\mathcal{C}^0_t\subseteq \mathcal{C}^\alpha_t$, $\mathcal{RS}^0_t\subseteq \mathcal{RS}^\alpha_t$, and for all $s\in \mathcal{RS}_t^0, \delta^0_t(s) = \delta^\alpha_t(s)$ and $\Theta^0_t(s)=\Theta^\alpha_t(s)$.


\subsubsection{Market State}

We additionally model the market as a discrete-time process $\mathcal{M}_\bullet := \{\mathcal{M}_t\}_{t\in\mathbb{N}}$.
At each timestep $t$, the market state is
\[
\mathcal{M}_t := \big(\{\mathcal{P}^\alpha_t\}_{\alpha\in\mathcal{N}},\;\mathcal{ML}_t\big)
\in \mathfrak{P}^n \times \mathfrak{L} = \mathfrak{M}\]
where $\mathcal{P}^\alpha_t\in\mathfrak{P}$ is agent $\alpha$'s private portfolio state and $\mathcal{ML}_t \in \mathfrak{L}$
is the public market ledger.

We assume a canonical balance/valuation map $\beta:\mathfrak{P}\to\mathbb{R}$ and define the balance of agent $\alpha$ at time $t$ by
\[
B^\alpha_t := \beta(\mathcal{P}^\alpha_t)\in\mathbb{R}.
\]



\subsubsection{Action Primitives}

We now specify the core action primitives of VAMP. These primitives will later be assembled into the
POSG transition and observation structure; for now, we define them as standalone kernels and update maps.

\paragraph{Discrete kernels.}
All stochastic primitives in this section have finite or countable output spaces, and may be treated as
probability mass functions (pmfs). Concretely, for an input space $X$ and discrete output space $Y$,
a stochastic kernel is a map $K:Y\times X\to[0,1]$ such that $\sum_{y\in Y}K(y\mid x)=1$ for all $x\in X$.

\paragraph{Proof Kernel}

Fix an agent $\alpha\in\mathcal{N}$. The proof kernel models allocating a time budget to attempt to resolve a target sequent.

% \paragraph{Type.}
Let the input space be
\[
X_{\mathrm{proof}} := \mathscr{L}\times\mathcal{S}\times \mathbb{N}^+,
\]
where $(s,\tau)$ consists of a target sequent $s=[\Gamma\proves\phi]$ and a discrete time budget $\tau$.
Let the output space be
\[
Y_{\mathrm{proof}} := \{(0,\bot)\}\;\cup\;\{(1,d): d\in\mathcal{P}_{fin}(\mathcal{F})\},
\]
where $(0,\bot)$ denotes failure, and $(1,d)$ denotes success together with a dependency set $d$.

The proof kernel for agent $\alpha$ is a pmf

\[\mathcal{K}^{\alpha}_{\mathrm{proof}} : Y_{\mathrm{proof}}\times X_{\mathrm{proof}} \to [0,1]\]


% \paragraph{Intended admissibility.}
Although $\mathcal{K}^{\alpha}_{\mathrm{proof}}$ is defined on all of $\mathcal{S}\times\mathbb{N}^+$,
in the intended game dynamics the agent is only allowed to invoke a proof attempt on \emph{concrete} targets
available in its local library (later captured via a state-dependent admissibility predicate).
On admissible inputs $(\mathcal{L},[\Gamma\proves\phi],\tau)$, a successful output $(1,d)$ is intended to satisfy $d\subseteq\Gamma$.

% \paragraph{Intended state effect (informal).}
Upon success, the environment will (i) mark the target sequent resolved in $\alpha$'s library and
(ii) record/update its dependency map at that sequent with the returned $d$, subject to global consistency
(e.g., preserving acyclicity of the induced dependency graph). Upon failure, no new resolved sequent is added.
These effects are enforced later by the global state-update operator.

\paragraph{Conjecture Kernel}

Fix an agent $\alpha\in\mathcal{N}$. The conjecture kernel models allocating a time budget to propose a new formula,
anchored at an existing sequent.

Let the input space be
\[
X_{\mathrm{conj}} :=\mathscr{L}\times \mathcal{S}\times \mathbb{N}^+,
\]
where $(s,\tau)$ consists of an anchor sequent $s=[\Gamma\proves\phi]$ and time budget $\tau$.
Let the output space be
\[
Y_{\mathrm{conj}} := \{(0,\bot)\}\;\cup\;\{(1,\psi): \psi\in\mathcal{F}\},
\]
where $(0,\bot)$ denotes failure and $(1,\psi)$ denotes successful conjecture of formula $\psi$.

The conjecture kernel for agent $\alpha$ is a pmf
\[
\mathcal{K}^{\alpha}_{\mathrm{conj}} : Y_{\mathrm{conj}}\times X_{\mathrm{conj}} \to [0,1].
\]

As with the proof kernel, $\mathcal{K}^{\alpha}_{\mathrm{conj}}$ is only intended to be invoked on anchors available to the agent (e.g., concrete sequents in $\alpha$'s local library).

Moreover, given admissible input $(\mathcal L, [\Gamma \proves \phi], \tau)$ and a successful output $(1,\psi) \sim \mathcal{K}_{\mathrm{conj}}^\alpha (\cdot \mid \mathcal L, [\Gamma \proves \phi], \tau)$, we enforce that $\psi$ is a ghost node. Namely, if $\mathcal{L} = (\mathcal{C},\mathcal{RS},\delta,\Theta)$, then $\psi \in \mathcal{F} \setminus \mathcal{C}$.


Upon such success with output $(1,\psi)$, we intend the environment to add $\psi$ to $\alpha$'s set of concrete formulas (and also add $\neg\psi$ to preserve the negation-pair invariant). This converts $\psi$ from ghost to concrete relative to $\alpha$.

\paragraph{Dependency Closure for Publication}

Publication is a deterministic primitive that copies a closure-completed subset of an agent's private library
into the shared library, and (by the shared-library invariant) into every agent's local library.

Fix agent $\alpha$ and time $t$, and suppose $\delta_t^\alpha$ induces the directed acyclic graph
$D(\delta_t^\alpha)=(\mathcal{CS}_t^\alpha,\mathcal{E}_t^\alpha)$ as defined earlier.
For any set of target resolved sequents $R\subseteq \mathcal{RS}_t^\alpha$, define its \emph{dependency closure}:
\begin{align*}
\mathrm{cl}^{\alpha}_t(R) := 
\Big\{\, &s\in \mathcal{RS}_t^\alpha : \exists r\in R \text{ such that there}\\& \text{is a directed path } s\to r
\text{ in } D(\delta_t^\alpha)\,\Big\}
\end{align*}

Equivalently, $\mathrm{cl}^\alpha_t(R)$ is the set of all resolved sequents required (transitively) as dependencies
of the sequents in $R$.


We now proceed to extend this closure operator to the payload for such publication actions. Namely, a publication payload is a pair $U_t^\alpha=(C,R)$ with $C\subseteq \mathcal{C}_t^\alpha,\;\; R\subseteq \mathcal{RS}_t^\alpha$.

Define its closure-completion $\overline{U_t^\alpha}:=\big(\overline{C},\overline{R}\big)$ where $\overline{R}:=\mathrm{cl}^\alpha_t(R)$ and
\[\overline{C}:= C \;\cup\; \bigcup_{[\Gamma\proves\phi]\in \overline{R}}(\Gamma\cup\{\phi\})\]

Thus, publishing resolved sequents forces publication of all formulas appearing in the antecedents and consequents of the dependency-closed witness set.


Therefore, given $\overline{U_t^\alpha}$, intuitively, the environment updates the shared library by unioning in $\overline{C}$ and $\overline{R}$ (and any associated dependency information it chooses to expose),
then propagates these additions to every agent's local library so that $\mathcal{L}^0_{t+1}\subseteq \mathcal{L}^\beta_{t+1}$
continues to hold for all $\beta$.



\paragraph{Query Model}

Each agent $\alpha$ maintains a stateful \emph{query model} which estimates how its prover stack will perform on a
requested proof search. Intuitively, the query model is a continually-updated predictor for
(i) the probability of successfully resolving a target sequent within a given time budget and
(ii) the expected time-to-solution (or a proxy thereof).


Fix an abstract space $\mathcal{Q}$ of query models (e.g., parameters of a Bayesian estimator, weights of a learned predictor, or a table of sufficient statistics). At each timestep $t$, agent $\alpha$ has a current query model $Q_t^\alpha \in \mathcal{Q}$.
A query operation takes as input a target (concrete) sequent:
\[
q^\alpha : \mathcal{Q}\times \mathcal{S} \to [0,1]\times \mathbb{N}
\]
and write
\[
(\hat p_t^\alpha(s),\hat \tau_t^\alpha(s))
:=
q^\alpha(Q_t^\alpha,s),
\]
where $\hat p_t^\alpha([\Gamma \proves \phi])$ is the model's estimate of $\Pr(\phi\text{ is true})$
and $\hat\tau_t^\alpha(s)$ is an estimate of the expected budget required to successfully prove $[\Gamma\proves \phi]$. 

More explicitly, for a fixed agent $\alpha$ and library state $\mathcal{L}$, and target sequent $s\in\mathcal{S}$, define the probability of success within budget $\tau$ by
\[
p^\alpha_\tau(s) \;:=\; \sum_{d\in\mathcal{P}_{fin}(\mathcal{F})}\mathcal{K}^{\alpha}_{\mathrm{proof}}\big((1,d)\mid \mathcal L, s,\tau\big).
\]
This induces a time-to-solution distribution by declaring a random variable $T^\alpha_s$ with CDF
\[
\Pr\!\big(T^\alpha_s \le \tau\big) \;:=\; p^\alpha_\tau(s), \qquad \tau\in\mathbb{N}^+.
\]
We interpret the eventual solvability (or ``truth/resolvability'' under $\alpha$'s prover stack) as
\[
p^\alpha_\infty(s) \;:=\; \sup_{\tau\in\mathbb{N}^+} p^\alpha_\tau(s),
\]
and, when $p^\alpha_\infty(s)>0$, the conditional expected solve time as
\[
\tau^\alpha_\infty(s) \;:=\; \mathbb{E}\!\left[T^\alpha_s \,\middle|\, T^\alpha_s < \infty\right]=\sum_{k\ge1}\Pr(T^\alpha_s\ge k \mid T^\alpha_s<\infty)
\]
Accordingly, the query readout $q^\alpha(Q^\alpha_t,s)=(\hat p^\alpha_t(s),\hat\tau^\alpha_t(s))$ is intended to approximate
\[
\hat p^\alpha_t(s)\approx p^\alpha_\infty(s),
\qquad
\hat\tau^\alpha_t(s)\approx \tau^\alpha_\infty(s),
\]
using the agent's accumulated evidence.



The query model is updated from observed outcomes of proof attempts (and optionally from public disclosures).

Define the space of query-training data as
\[
\mathcal{D}(\mathcal{Q}) := \mathcal{S} \times (\mathbb{N}^+ \times \mathcal{N}),
\]
where a datum $(s,(\tau, \alpha'))$ denotes a successful proof attempt by agent $\alpha'$ on $s$ on time budget $\tau$.

Define an update operator
\[
\Delta^\alpha : \mathcal{Q}\times \mathcal{D}(\mathcal{Q}) \to \mathcal{Q}
\]

Given an observed training datum $d_t^\alpha\in\mathcal{D}(\mathcal{Q})$ available to agent $\alpha$ at time $t$, the query model updates as
\[
Q_{t+1}^\alpha :=\Delta^\alpha(Q_t^\alpha, d_t^\alpha),
\]
We enforce the commutativity and associativity of $\Delta^\alpha$, so that if multiple data points $\{(d_t^\alpha)_i\}_{i=0}^{k-1}$ are observed in a single timestep, there is a single and well-defined batch update which is the composition (in any order) of the individual datum updates: i.e., if $Q_t^\alpha=(Q_t^\alpha)_0$,  
\[
(Q_{t}^\alpha)_{i+1} = \Delta^\alpha((Q_t^\alpha)_i, (d_t^\alpha)_i)
\]
\[Q_{t+1}^\alpha := (Q_{t}^\alpha)_{k-1}\]
which is consistent across all possible permutations on $\{0,\cdots,k-1\}$.



The framework does not prescribe the internal form of $Q_t^\alpha$, and allows downstream experiments to arbitrarily instantiate $(\mathcal{Q}, q^\alpha, \Delta^\alpha)$ as desired, whether that means querying based on Bayesian update rules, learned predictors, or any other mechanism.

\paragraph{Market Mechanism}
\label{sec:vamp-market}

We briefly take an aside to define the \emph{market mechanism} as a modular plugin that that governs (i) which market actions are available, (ii) which are admissible at a given world state, and (iii) how the market state evolves. One may parameterize and study many different verifiable allocation games by simply switching out this underlying economic reward mechanism.

We define a market mechanism $\mathscr{M}$ as a tuple:
\begin{align*}
\mathscr{M} :=\Big(&\mathcal{N}, \mathfrak{M},\beta, \mathcal{W},\{\mathcal{A}^\alpha_{\mathrm{mkt}}\}_{\alpha\in\mathcal{N}},
\\&\{\mathrm{Adm}^\alpha_{\mathrm{mkt}}\}_{\alpha\in\mathcal{N}},\mathcal{T}_{\mathrm{mkt}}\Big)
\end{align*}
where:
\begin{itemize}
    \item $\mathcal{N}$ is the set of agents, indexed by $\alpha\in \mathcal{N}$.
    \item $\mathfrak{M}=(\mathfrak{P}^n,\mathfrak{L})$ is the market state space, where $\mathfrak{P}$ is the space of possible portfolios and $\mathfrak{L}$ is the space of possible public market ledgers.
    \item $\beta : \mathfrak{P} \to \mathbb{R}$ is the canonical balance/valuation map.
    \item $\mathcal{W}$ is the world (state) space. Note that $\mathfrak{M}$ is a component in $\mathcal{W}$.
  
    \item $\mathcal{A}^\alpha_{\mathrm{mkt}}$ is the (ambient) \textit{market} action space of agent $\alpha$. Note that the market no-op is always a valid action: $\bot_{\mathrm{mkt}} \in \mathcal{A}_{\mathrm{mkt}}^\alpha$. The joint market action space is:
    \[
\mathcal{A}_{\mathrm{mkt}} := \prod_{\alpha\in\mathcal{N}} \mathcal{A}^\alpha_{\mathrm{mkt}}
\]
    \item $\mathrm{Adm}^\alpha_{\mathrm{mkt}}:\mathcal{W}\times\mathcal{A}^\alpha_{\mathrm{mkt}}\to\{0,1\}$ is the market admissibility predicate for agent $\alpha$. It induces the state-dependent admissible market action set:
    \[
    \mathcal{A}^\alpha_{\mathrm{mkt}}(s) := \{a^\alpha\in\mathcal{A}_{\mathrm{mkt}}^\alpha : \mathrm{Adm}_{\mathrm{mkt}}^\alpha(s,a^\alpha)=1\}.
    \]
    \item $\mathcal{T}_{\mathrm{mkt}} : \Big(\mathcal{W} \times \mathcal{W} \times \mathcal{A}_{\mathrm{mkt}}\Big) \times \Sigma_{\mathcal{W}} \to [0,1]$ is a stochastic transition kernel with input space $\mathcal{W} \times \mathcal{W} \times \mathcal{A}_{\mathrm{mkt}}$ and output space $\mathcal{W}$. Namely, given a change in world state from $w \to w' \in \mathcal{W}$, and a joint market action $a_{\mathrm{mkt}}\in\mathcal{A}_{\mathrm{mkt}}$, this kernel induces a distribution over the next world state:
\[
\bar{w}\sim \mathcal{T}_{\mathrm{mkt}}(\cdot\mid w,w',a_{\mathrm{mkt}}).
\]
Deterministic mechanisms are recovered as Dirac kernels.

\end{itemize}




\subsection{Verifiable Allocation Market POSG}


We now define the \emph{Verifiable Allocation Market POSG} (VAMP) as a particular POSG designed to model multiagent verifiable subtask allocation mediated by an arbitrary market mechanism -- interpreted as an abstracted decentralized formal theorem proving environment.

Fix a set of agents $\mathcal{N}$. A VAMP is a class of infinite-horizon discounted POSG's, parameterized by a particular market mechanism $\mathscr{M}$:
\[
\Big(\mathcal{N},\mathcal{W},\{\mathcal{O}^\alpha\}_{\alpha\in\mathcal{N}},\{\mathcal{Z}^\alpha\}_{\alpha\in\mathcal{N}},\{\mathcal{A}^\alpha\}_{\alpha\in\mathcal{N}},
\]\[\{\mathrm{Adm}^\alpha\}_{\alpha\in\mathcal{N}},\mathcal{T},
\{r^\alpha\}_{\alpha\in\mathcal{N}},\gamma,b_0, \mathscr{M}\Big)
\]
defined as follows.

\subsubsection{World State Space $\mathcal{W}$}

A world state at time $t$ is a tuple
\begin{align*}
w_t :=
\Big(
&\{\mathcal{L}^\alpha_t\}_{\alpha\in\mathcal{N}},
\ \mathcal{L}^0_t,
\ \mathcal{M}_t,
\ \{J^\alpha_t\}_{\alpha\in\mathcal{N}},
\ \{Q^\alpha_t\}_{\alpha\in\mathcal{N}},
\\&\ \{\Omega^\alpha_t\}_{\alpha\in\mathcal{N}},
\ \{U_{t,\mathrm{proof}}^\alpha\}_{\alpha\in\mathcal N},
\ \{U_{t,\mathrm{conj}}^\alpha\}_{\alpha\in\mathcal N}
\Big)\in\mathcal{W},
\end{align*}
where:
\begin{itemize}
\item $\mathcal{L}^\alpha_t$ is agent $\alpha$'s private library and
$\mathcal{L}^0_t$ is the shared public library.
\item $\mathcal{M}_t=(\{\mathcal{P}^\alpha_t\}_{\alpha\in\mathcal{N}},\;\mathcal{ML}_t) \in \mathfrak{P}^n \times \mathfrak{L} =\mathfrak{M}$ is the market state, as defined in $\mathscr{M}$.
\item $J^\alpha_t$ is the durative job descriptor
\[
J^\alpha_t \in \{\bot\}\ \cup\ \big(\{\mathrm{Proof}, \mathrm{Conj}\}\times\mathcal{S}\times\mathbb{N}^+\times\mathbb{N}^+\big)
\]
with $\bot$ meaning idle and $(\cdot,s,\tau_{\mathrm{rem}},\tau_{\mathrm{eff}})$ meaning:
(i) $\tau_{\mathrm{rem}}$ steps remain in the \emph{current run} and
(ii) $\tau_{\mathrm{eff}}$ is the \emph{effective cumulative budget} used to evaluate the run upon completion.
Concretely, when a run is initiated at time $t$ with intended additional duration $\tau_{\mathrm{dur}}$ on target $s$,
we set $\tau_{\mathrm{eff}} := U^\alpha_{t,\mathrm{-}}(s) + \tau_{\mathrm{dur}}$ and $\tau_{\mathrm{rem}} := \tau_{\mathrm{dur}}$.
The agent is forced to take $\tau_{\mathrm{dur}}$ consecutive no-op actions, and at completion the outcome is sampled from the
appropriate kernel using budget $\tau_{\mathrm{eff}}$.
\item $Q^\alpha_t\in\mathcal{Q}$ is agent $\alpha$'s query model state, updated by $\Delta^\alpha$ on successful proofs.
\item  $\Omega^\alpha_t \in \{\bot\}\ \sqcup\ (\mathcal{S}\times[0,1]\times\mathbb{N})$ is agent $\alpha$'s last query response, either empty $\bot$ or a triple $(s,\hat p,\hat \tau)$ signifying an estimate that the sequent $s$ is true with probability $\hat p$, and given it is true, an estimate of the time it would take to prove it $\hat \tau$.
\item $U^\alpha_{t,\mathrm{-}} : \mathcal S \to \mathbb N$ records the total time already spent by agent $\alpha \in \mathcal N$ attempting a proof or conjecture for sequent $s \in \mathcal S$, up to time $t$. This is used to allow agents to ``pick up where they left off''. Note $U^\alpha_{0,\mathrm{-}}\equiv 0$.
\end{itemize}
We take $\mathcal{W}$ to be the set of tuples satisfying all such tuples satisfying the consitituent objects' various invariants and conditions.


\subsubsection{Observations $\mathcal{O}^\alpha$ and Observation Kernels $\mathcal{Z}^\alpha$}

% Public events are defined as the change in the public projection of the world state:
% \[
% \mathrm{Pub}(w_t):=(\mathcal{L}^0_t,\Theta^0_t,\mathcal{ML}_t),
% \]\[E^{\mathrm{pub}}_t := \mathrm{Pub}(w_{t+1}) \setminus \mathrm{Pub}(w_{t}),
% \]
% Thus, any update to a public field is automatically publicly visible.

Agent $\alpha$ observes
\[
o^\alpha_t :=
\big(\mathcal{L}^\alpha_t,\mathcal{P}^\alpha_t,B^\alpha_t, \Omega_t^\alpha;\ \mathcal{L}^0_t,\mathcal{ML}_t\big)%;\ E^{\mathrm{pub}}_t\big)
\in\mathcal{O}^\alpha.
\]
The observation kernel is deterministic:
\[
\mathcal{Z}^\alpha(o^\alpha\mid w,a,w')=\mathbf{1}[\,o^\alpha=\mathrm{Obs}^\alpha(w)\,],
\]
where $\mathrm{Obs}^\alpha(w)$ computes the observation slice as above.


\subsubsection{Action Space $\mathcal{A}^\alpha$ and Admissibility $\mathrm{Adm}^\alpha$}

For each agent $\alpha$, define the tagged \textit{non-market} action space
\begin{align*}
\mathcal{A}_{\mathrm{nm}}^\alpha
:=
&
\{\mathrm{Prove}\}\times(\mathcal{S}\times\mathbb{N}^+)
\\&\sqcup
\{\mathrm{Conj}\}\times(\mathcal{S}\times\mathbb{N}^+)
\\&\sqcup
\{\mathrm{Pub}\}\times(\mathcal{P}_{fin}(\mathcal{F})\times \mathcal{P}_{fin}(\mathcal{S}))
\\&\sqcup
\{\mathrm{Qry}\}\times \mathcal{S}
\\&\sqcup\{\bot_{\mathrm{nm}}\}
\end{align*}
Where $\bot_{\mathrm{nm}}$ indicates a non-market no-op. Thus, the full action space for agent $\alpha$ may be written as:
\[
\mathcal{A}^\alpha = \mathcal{A}^\alpha_{\mathrm{nm}} \sqcup \mathcal{A}^\alpha_{\mathrm{mkt}}
\]
where $\mathcal{A}^\alpha_{\mathrm{mkt}}$ is specified in $\mathscr{M}$. Note then that as an agent $\alpha$ may only take one action at a given timestep, both $\bot_{\mathrm{nm}}$ and $\bot_{\mathrm{mkt}}$ would indicate a no-op.

As such, the joint action space is 
\[\mathcal{A}:=\prod_{\alpha\in\mathcal{N}}\mathcal{A}^\alpha\]


Note that for a  given joint action $a \in \mathcal{A}$, we may uniquely decompose it into a non-market and market joint action. In other words:
\[a = a_{\mathrm{nm}} \oplus a_{\mathrm{mkt}}\]
Where for each agent-component of the joint action $a^\alpha \in \mathcal{A}^\alpha$, if $a^\alpha \in \mathcal{A}^\alpha_{\mathrm{nm}}$, then let $a^\alpha_{\mathrm{nm}} = a^\alpha$ and $a^\alpha_{\mathrm{mkt}} = \bot_{\mathrm{mkt}}$. Alternatively, if $a^\alpha \in \mathcal{A}^\alpha_{\mathrm{mkt}}$, then let $a^\alpha_{\mathrm{mkt}} = a^\alpha$ and $a^\alpha_{\mathrm{nm}} = \bot_{\mathrm{nm}}$. Thus, we may define:
\[a_{\mathrm{nm}} = \prod_{\alpha \in \mathcal{N}} a^\alpha_{\mathrm{nm}},\qquad a_{\mathrm{mkt}} = \prod_{\alpha \in \mathcal{N}} a^\alpha_{\mathrm{mkt}}\]




Now, define $\mathrm{Adm}^\alpha:\mathcal{W}\times\mathcal{A}^\alpha\to\{0,1\}$ as an admissibility predicate, which enforces semantic and operational constraints. 

Fix $w_t \in \mathcal{W}$ and consider $\alpha \in \mathcal{N}$. Then:
\begin{itemize}
    \item If $J^\alpha_t=\bot$, then:
\begin{itemize}
\item $\mathrm{Adm}^\alpha(w_t,\bot)=1$.
\item For $a^\alpha_t = (\mathrm{Prove},(s,\tau))$ and $a^\alpha_t = (\mathrm{Conj},(s,\tau))$:
\[
\mathrm{Adm}^\alpha(w_t,a^\alpha_t)=1 \iff s \in \mathcal{CS}_t^\alpha
\]
\item For $a^\alpha_t = (\mathrm{Pub},(C,R))$: \[\mathrm{Adm}^\alpha(w_t,a^\alpha_t)=1 \iff C\subseteq\mathcal{C}_t^\alpha,\ R\subseteq\mathcal{RS}_t^\alpha\]
\item For $a^\alpha_t = (\mathrm{Qry},s)$: 
\[\mathrm{Adm}^\alpha(w_t,a_t^\alpha)=1 \iff s \in \mathcal{CS}_t^\alpha\]
\item For market actions $a^\alpha_t\in\mathcal{A}^\alpha_{\mathrm{mkt}}$:
\[
\mathrm{Adm}^\alpha(w_t,a^\alpha_t)= \mathrm{Adm}^\alpha_{\mathrm{mkt}}(w_t,a^\alpha_t) 
\]
where $\mathrm{Adm}_{\mathrm{mkt}}^\alpha$ is specified by $\mathscr{M}$.
\end{itemize}

\item If $J^\alpha\neq\bot$, then the agent is forced to take no-op: $\mathrm{Adm}^\alpha(w_t,\bot)=1$ and for all $a^\alpha_t\in \mathcal{A}^\alpha\setminus\{\bot\}$, $\mathrm{Adm}^\alpha(w_t,a^\alpha_t)=0$.
\end{itemize}

\subsubsection{Transition Kernel $\mathcal{T}$}

We define $\mathcal{T}(\cdot\mid w_t,a_t)$ by a two-stage generative construction. Namely, given an admissible input pair $(w_t,a_t) \in \mathcal{W} \times \mathcal{A}$, uniquely decompose the joint action as
\[
a_t = a_{t,\mathrm{nm}} \oplus a_{t,\mathrm{mkt}}
\]

We then define a non-market transisition kernel $\mathcal{T}_{\mathrm{nm}} : (\mathcal{W} \times \mathcal{A}_\mathrm{nm}) \times \Sigma_{\mathcal{W}} \to [0,1]$ and consider the market transition kernel $\mathcal{T}_{\mathrm{mkt}} : (\mathcal{W} \times \mathcal{W} \times \mathcal{A}_\mathrm{mkt}) \times \Sigma_{\mathcal{W}} \to [0,1]$ as specified by $\mathscr{M}$. Intuitively, the full transition kernel $\mathcal{T}$ is the composition of the two stages:
first draw $w'_t\sim \mathcal{T}_{\mathrm{nm}}(\cdot\mid w_t,a_{t,\mathrm{nm}})$, then draw
$w_{t+1}\sim \mathcal{T}_{\mathrm{mkt}}(\cdot\mid w_t,w'_t,a_{t,\mathrm{mkt}})$.

We now proceed to define each stage of this construction.

\paragraph{Background: Deterministic Pre-Update}
 
 We introduce a deterministic map that serves as the first step of the non-market transition kernel. It applies everything that happens immediately from the joint non-market action at this step, before any stochastic proof/conjecture outcomes are realized. So, given a world state $w \in \mathcal{W}$, we define:
\[\tilde w := \mathrm{Pre}(w,\mathbf a_{\mathrm{nm}})\]
as the world state obtained by applying the following deterministic edits in order:
\begin{enumerate}
    \item Decrement job timers. Namely, for each $\alpha \in \mathcal{N}$:
    \begin{itemize}
        \item If $J^\alpha=\bot$, leave it.
	\item If $J^\alpha=(\mathrm{Proof},s,\tau_{\mathrm{rem}},\tau_{\mathrm{tot}})$ or $(\mathrm{Conj},s,\tau_{\mathrm{rem}},\tau_{\mathrm{tot}})$ with $\tau_{\mathrm{rem}}>0$, set
$\tau_{\mathrm{rem}}\leftarrow \tau_{\mathrm{rem}}-1$.
    \end{itemize}
\item Start new proof/conjecture jobs if idle. Namely, for each $\alpha \in \mathcal N$ with $J^\alpha=\bot$:
\begin{itemize}
    \item 	If $a^\alpha_{\mathrm{nm}}=(\mathrm{Prove},(s,\tau))$, set
$J^\alpha \leftarrow (\mathrm{Proof},s,\tau,\tau+U^\alpha_{\mathrm{proof}}(s))$.
	\item	If $a^\alpha_{\mathrm{nm}}=(\mathrm{Conj},(s,\tau))$, set
$J^\alpha \leftarrow (\mathrm{Conj},s,\tau,\tau+U^\alpha_{\mathrm{conj}}(s))$.
	\item	Otherwise do nothing.
\end{itemize}
\item Apply publication actions and propagate to shared library. Namely, for each $\alpha\in\mathcal{N}$, if $a^\alpha_{\mathrm{nm}}=(\mathrm{Pub},(C^\alpha,R^\alpha))$, then consider the payload $U^\alpha:=(C^\alpha,R^\alpha)$. Compute the closure $\overline{U^\alpha}$ and update the shared library and query model. Namely, let $\mathcal{C}^0_\mathrm{old} := \mathcal{C}^0$ and $\mathcal{RS}^0_\mathrm{old} := \mathcal{RS}^0$, and update:
\[\mathcal C^0 \leftarrow \mathcal C^0 \cup \bigcup_{\alpha}\overline C^\alpha,
\qquad
\mathcal{RS}^0 \leftarrow \mathcal{RS}^0 \cup \bigcup_{\alpha}\overline R^\alpha\]

Then for all $s \in \mathcal{RS}^0 \setminus \mathcal{RS}^0_\mathrm{old}$:
\[\delta^0(s) \leftarrow \delta^\alpha(s),\qquad  \Theta^0(s) \leftarrow \Theta^\alpha(s)\]

Then propagate public additions to each private library to maintain the invariant $\mathcal L^0\subseteq \mathcal L^\beta$:
\[\mathcal C^\beta \leftarrow \mathcal C^\beta \cup (\mathcal C^0\setminus \mathcal C^0_{\text{old}})\]\[
\mathcal{RS}^\beta \leftarrow \mathcal{RS}^\beta \cup (\mathcal{RS}^0\setminus \mathcal{RS}^0_{\text{old}})\]
Also update for all $s\in \mathcal{RS}^0 \setminus \mathcal{RS}^0_\mathrm{old}$:
 \[\delta^\beta(s)\leftarrow\delta^0(s),\qquad \Theta^\beta(s)\leftarrow\Theta^0(s)\]
 \[Q^\beta \leftarrow \Delta^\beta(Q^\beta,(s,\Theta^0(s)))\]
 


\item Update query buffer. Namely, for each $\alpha\in\mathcal{N}$, if $a^\alpha_{\mathrm{nm}}=(\mathrm{Qry},s)$, set
\[\Omega^\alpha \leftarrow \big(s,\hat p,\max\{0,\hat\tau-U^\alpha_{\mathrm{proof}}(s)\}\big)\]
where $(\hat p,\hat\tau):=q^\alpha(Q^\alpha,s)$. Otherwise set $\Omega^\alpha\leftarrow \bot$.

\end{enumerate}

\paragraph{Background - Stochastic Job Completion}

Given $\tilde w$, define completion index sets:
\[
    \mathsf{Comp}_{\mathrm{proof}}(\tilde w) = \{\alpha : J^\alpha(\tilde w)=(\mathrm{Proof},s,0,\tau)\text{ for some }s,\tau\}
\]
\[\mathsf{Comp}_{\mathrm{conj}}(\tilde w) = \{\alpha : J^\alpha(\tilde w)=(\mathrm{Conj},s,0,\tau)\text{ for some }s,\tau\}\]

For each $\alpha\in \mathsf{Comp}_{\mathrm{proof}}(\tilde w)$, let $(s^\alpha_{\mathrm{proof}}(\tilde w),\tau^\alpha_{\mathrm{proof}}(\tilde w))$ be the $(s,\tau)$ stored in $J^\alpha(\tilde w)$ (and similarly for conjecture completions).

Then, we may now define outcome spaces (per completing job):
\[Y^{\alpha}_{\mathrm{proof}} := \{(0,\bot)\}\cup\{(1,d): d\in\mathcal P_{\mathrm{fin}}(\mathcal F)\}\]
\[Y^{\alpha}_{\mathrm{conj}} := \{(0,\bot)\}\cup\{(1,\psi): \psi\in\mathcal F\}\]
And the joint outcome space, dependent on $\tilde w$:
\[Y(\tilde w)
:=
\prod_{\alpha\in\mathsf{Comp}_{\mathrm{proof}}(\tilde w)} Y^{\alpha}_{\mathrm{proof}}
\;\times\;
\prod_{\alpha\in\mathsf{Comp}_{\mathrm{conj}}(\tilde w)} Y^{\alpha}_{\mathrm{conj}}\]

Thus, we may define a product measure over this joint outcome space by conditional independence across completing jobs:
\begin{align*}
&\mathbb K_{\tilde w}(dy)
:=
\\&\left(\prod_{\alpha\in\mathsf{Comp}_{\mathrm{proof}}(\tilde w)}
\mathcal K^\alpha_{\mathrm{proof}}(dy^\alpha \mid \mathcal L^\alpha(\tilde w), s^\alpha_{\mathrm{proof}}(\tilde w),\tau^\alpha_{\mathrm{proof}}(\tilde w))
\right)
\\&\cdot\left(\prod_{\alpha\in\mathsf{Comp}_{\mathrm{conj}}(\tilde w)}
\mathcal K^\alpha_{\mathrm{conj}}(dy^\alpha \mid \mathcal L^\alpha(\tilde w), s^\alpha_{\mathrm{conj}}(\tilde w),\tau^\alpha_{\mathrm{conj}}(\tilde w))
\right)
\end{align*}

\paragraph{Background - Deterministic Post-Update} 

We now aim to deterministically apply the realized outcomes $y\in Y(\tilde w)$ to update libraries, dependency maps, solve-times, query models, and clear completed jobs.

Given $\tilde w$ and $y\in Y(\tilde w)$, define
\[w^+ := \mathrm{Post}_{\mathrm{nm}}(\tilde w,y)\]
by applying:
\begin{enumerate}
    \item Resolve completing proof jobs. Namely, for each $\alpha\in\mathsf{Comp}_{\mathrm{proof}}(\tilde w)$, let
\[y^\alpha = (b^\alpha, z^\alpha)\in Y^\alpha_{\mathrm{proof}}\]
Let $(s^\alpha,\tau^\alpha)$ be the job’s stored target and total budget. Regardless of proof success, set $U^\alpha_{\mathrm{proof}}(s^\alpha)\ \leftarrow\ \tau^\alpha$.

Then, if $b^\alpha=1$ and $z^\alpha=d$, then update:
\begin{itemize}
    \item	$\mathcal{RS}^\alpha \leftarrow \mathcal{RS}^\alpha \cup \{s^\alpha\}$
	\item	$\delta^\alpha(s^\alpha) \leftarrow d$ (after enforcing $d\subseteq \Gamma$ / acyclicity checks)
	\item	$\Theta^\alpha(s^\alpha)\leftarrow (\tau^\alpha,\alpha)$
	\item	$Q^\alpha \leftarrow \Delta^\alpha(Q^\alpha,(s^\alpha,\Theta^\alpha(s^\alpha)))$
\end{itemize}
	\item Resolve completing conjecture jobs. Namely, for each $\alpha\in\mathsf{Comp}_{\mathrm{conj}}(\tilde w)$, let
\[y^\alpha = (b^\alpha, z^\alpha)\in Y^\alpha_{\mathrm{conj}}\]
Regardless of conjecture success, set $U^\alpha_{\mathrm{conj}}(s^\alpha)\ \leftarrow\tau^\alpha$. Then, if $b^\alpha=1$ and $z^\alpha=\psi$, then:
\[\mathcal C^\alpha \leftarrow \mathcal C^\alpha \cup \{\psi,\neg\psi\}\]
\item Clear completed jobs. Namely, for each $\alpha\in \mathsf{Comp}_{\mathrm{proof}}(\tilde w)\cup \mathsf{Comp}_{\mathrm{conj}}(\tilde w)$, set:
$J^\alpha \leftarrow \bot$.

\end{enumerate}



\paragraph{Stage I: Non-market transition $\mathcal{T}_{\mathrm{nm}}$}

The non-market transition updates all components of the world state except the market state $\mathcal{M}_t$.
Formally, we define a kernel
\[
\mathcal{T}_{\mathrm{nm}}:\ (\mathcal{W}\times \mathcal{A}_{\mathrm{nm}})\times \Sigma_{\mathcal{W}}\to[0,1]
\]
which is implemented as the following pushforward kernel:
\[\mathcal T_{\mathrm{nm}}(U\mid w,a_{\mathrm{nm}})
\;=\;
\int_{Y(\widetilde w)}
\mathbf 1\!\left[\mathrm{Post}_{\mathrm{nm}}(\widetilde w,y)\in U\right]\,
\mathbb K_{\widetilde w}(dy)\]
For measurable $U\subseteq \mathcal W$, $\widetilde w:=\mathrm{Pre}(w,a_{\mathrm{nm}})$.


\paragraph{Stage II: Market transition via $\mathscr{M}$}

The market mechanism $\mathscr{M}$ supplies the market transition kernel $\mathcal{T}_{\mathrm{mkt}}$.
Given the realized non-market update $w'_t \sim \mathcal{T}_\mathrm{nm}(\cdot \mid w_t, a_{t,\mathrm{nm}})$ and the market joint action $a_{t,\mathrm{mkt}}$,
we sample the final next state by
\[
w_{t+1}\sim \mathcal{T}_{\mathrm{mkt}}(\cdot\mid w_t,w'_t,a_{t,\mathrm{mkt}}).
\]
Intuitively, $\mathcal{T}_{\mathrm{mkt}}$ may update portfolios $\{\mathcal{P}^\alpha\}$ and the public ledger $\mathcal{ML}$
as a function of both the submitted market actions and the realized mathematical progress encoded in the state change $w_t\to w'_t$,
e.g.\ settlement on newly published proofs, bounty claims, liquidations, AMM price updates, etc. Of course, however, the actual implementation of $\mathcal{T}_\mathrm{mkt}$ is entirely dependent on the specification of the market mechanism $\mathscr{M}$.

\paragraph{Full transition kernel}


Finally, we may define, for any measurable $U\subseteq\mathcal{W}$,
\begin{align*}
\mathcal{T}(U\mid w_t,a_t) = \int_{\mathcal{W}}
\mathcal{T}_{\mathrm{mkt}}(U\mid w_t,w',a_{t,\mathrm{mkt}})
\ \mathcal{T}_{\mathrm{nm}}(dw'\mid w_t,a_{t,\mathrm{nm}})
\end{align*}



\subsubsection{Reward, Discount Factor, and Initial Distribution}

We take the per-step reward to be the (expected) \emph{change in balance}:
\[
r^\alpha(w,a) := \mathbb{E}[B^\alpha(w') - B^\alpha(w) \mid w' \sim \mathcal T(\cdot \mid w,a)]
\]

Where $B^\alpha : \mathcal{W} \to \mathbb{R}$ is given by $B^\alpha(w) = \beta(P^\alpha(w))$, where $P^\alpha(w)$ is the usual projection of agent $\alpha$'s portfolio from the $\mathcal{M}$ component of $w \in \mathcal{W}$ and $\beta$ is defined by the $\mathscr{M}$ component of $w \in \mathcal{W}$.

% \subsubsection{Initial Distribution}

Finally, $b_0$ is a distribution over initial world states $w_0\sim b_0$ specifying initial private/shared libraries, initial job statuses
(typically idle), initial solve-time maps (typically empty), initial query models $Q^\alpha_0$, and the initial market state, supported on
states satisfying all invariants.

\subsubsection{Consistency and well-definedness}

We note that we require a VAMP to be \textit{consistent} with its market mechanism $\mathscr{M}$, meaning that if we have a market mechanism:
\begin{align*}
\mathscr{M} =\Big(&\mathcal{N}, \mathfrak{M},\beta, \mathcal{W},\{\mathcal{A}^\alpha_{\mathrm{mkt}}\}_{\alpha\in\mathcal{N}},
\\&\{\mathrm{Adm}^\alpha_{\mathrm{mkt}}\}_{\alpha\in\mathcal{N}},\mathcal{T}_{\mathrm{mkt}}\Big)
\end{align*}

The following tuple:
\[
\Big(\mathcal{N}',\mathcal{W}',\{\mathcal{O}^\alpha\}_{\alpha\in\mathcal{N}},\{\mathcal{Z}^\alpha\}_{\alpha\in\mathcal{N}},\{\mathcal{A}^\alpha\}_{\alpha\in\mathcal{N}},
\]\[\{\mathrm{Adm}^\alpha\}_{\alpha\in\mathcal{N}},\mathcal{T},
\{r^\alpha\}_{\alpha\in\mathcal{N}},\gamma,b_0, \mathscr{M}\Big)
\]
is a VAMP if and only if $\mathcal N' = \mathcal N, \mathcal W' = \mathcal W$, $\mathcal{A}^\alpha_\mathrm{mkt}\subseteq\mathcal{A}^\alpha$, for $a \in \mathcal{A}^\alpha_\mathrm{mkt}$ and $w \in \mathcal W$,$\mathrm{Adm}^\alpha_\mathrm{mkt}(w,a)=\mathrm{Adm}^\alpha(w,a)$. 

In other words, we may consider the following generator:
\[
\Big(
\mathcal{N},\mathfrak{P},\mathfrak{L}, \beta, \{\mathcal{A}^\alpha_{\mathrm{mkt}}\}_{\alpha\in\mathcal{N}},\{\mathrm{Adm}^\alpha_{\mathrm{mkt}}\}_{\alpha\in\mathcal{N}},\mathcal{T}_{\mathrm{mkt}},\gamma,b_0
\Big)
\]
Which uniquely specifies a well-defined VAMP using the above formulation.


\subsection{Equilibria and Tractibility}

[todo: study solution concepts, existence of equilibria, and tractability of computing them]

\subsection{Comparison to Centralized Mechanisms}


[todo: prove PoA-like results as mentioned earlier, showing that VAMP is better than centralized mechanisms (in cenrtain specializations) in private info settings. and conjecture that this is generally true too?]